\section{Equality Characterization}\label{sec:equality}
%==============================================================================

\medskip
\noindent\textbf{Classification of the non-clique boundary.} The remaining problem in this section is to characterize which upward‑closed agreement‑set families produce transitive confusability, and therefore trigger the cluster‑graph equality collapse. This does not weaken the equality route established here. Trivial single-view families and the fiber-coherent/meet-witnessed families already give broad sufficient conditions under which the collapse is completely determined; these are precisely the dominant patterns in architectures governed by a canonical transcript fiber or a common refining subview. The unresolved part is the residual non-clique boundary beyond those structural patterns. The binary square remains the smallest witness that this residual boundary is nonempty. A plausible conjectural direction is that transitivity of $G_{\mathrm{conf}}$ is governed by intersection-closure properties of the underlying admissible view family $\mathcal V$.

The original clique-fiber regime is now isolated as an equality subclass rather than only a threshold corollary. If a graph is generated by a surjective label map $x \mapsto b(x)$ with adjacency exactly between distinct vertices in the same fiber, then one representative from each fiber forms an independent set of size $|\mathcal B|$, while the complement graph is colorable by the same label map using $|\mathcal B|$ colors. This equality for cluster-graph structure is classical; the mechanized development packages the model-specific export needed here into a restricted-class theorem:
\[
\Theta(G)=\vartheta_{\infty}(G)=\log |\mathcal B|.
\]
Transitivity of confusability exports the model-generated graph directly to the classical cluster-graph setting through the following mechanized implication:
\[
\text{confusability is transitive}
\Longrightarrow
G_{\mathrm{conf}}=\mathrm{Cluster}(\Pi_{\mathrm{cc}}).
\]
Hence, if confusability is transitive, the base confusability graph is exactly the cluster graph on connected components, and the model-specific equality theorem yields
\[
\Theta(G_{\mathcal V})=\vartheta_\infty(G_{\mathcal V})=\log |\Pi_{\mathrm{cc}}|,
\]
where $\Pi_{\mathrm{cc}}$ is the connected-component partition of the base confusability graph. Between arbitrary base models and full fiber coherence, a checkable sufficient condition is that the view family be \emph{meet-witnessed}: every pair of allowed views contains another allowed view inside their intersection. Under that condition, any two confusability witnesses can be pulled back to a common subview, so confusability is transitive and the same equality follows. Fiber coherence is then a stronger sufficient structural condition: if equality at one allowed location already fixes the full observation transcript, transitivity follows, the connected components are exactly the realized transcript fibers, and the equality sharpens further to
\[
\Theta(G_{\mathcal V})=\vartheta_\infty(G_{\mathcal V})=\log |\mathcal T_{\mathrm{real}}|.
\]
At the witness level, this same collapse admits a local composition form: base confusability is the edge union of the single-view cluster graphs, and transitivity is equivalent to closure of those fixed-view witnesses under two-step composition through an intermediate state. This is the exact local criterion for the current equality mechanism: the collapse to a cluster graph is determined by how single-view witnesses compose, not by an opaque global coincidence of the full confusability graph. Thus the equality question is structural: it identifies when genuinely non-clique failure geometry collapses to exact cluster behavior.\leanmeta{\LHrng{MFT}{85}{91}, \LHrng{MFT}{100}{102}, \LHrng{MFT}{107}{109}, \LHrng{MFT}{136}{137}}

\paragraph{Running example: Binary square.}
The binary square of Section~\ref{sec:binary-square} does \emph{not} have transitive confusability: $(0,0) \sim (0,1)$ and $(0,1) \sim (1,1)$, but $(0,0) \not\sim (1,1)$. The 4-cycle is not a cluster graph, so the transitivity-based equality collapse does not apply. In this special case the classical graph invariants still satisfy $\Theta(C_4)=\vartheta_\infty(C_4)=\log 2$, but that coincidence no longer comes from the cluster-collapse mechanism isolated in this section. The binary square therefore marks the boundary of the paper's equality route: it is the first non-clique witness, yet not a transitive one.

\medskip
\noindent\textbf{Hierarchy summary.} Table~\ref{tab:equality-hierarchy} summarizes the hierarchy of sufficient conditions for the cluster-graph collapse.

\begin{table}[t]
\centering
\begingroup
\small
\setlength{\tabcolsep}{4pt}
\renewcommand{\arraystretch}{1.1}
\renewcommand{\tabularxcolumn}[1]{m{#1}}
\begin{tabularx}{0.98\linewidth}{@{}>{\RaggedRight\arraybackslash}m{0.21\linewidth}>{\RaggedRight\arraybackslash}m{0.16\linewidth}>{\RaggedRight\arraybackslash}m{0.31\linewidth}>{\RaggedRight\arraybackslash}m{0.22\linewidth}@{}}
\toprule
\textbf{Condition} & \textbf{Logical status} & \textbf{Structural consequence} & \textbf{Capacity value} \\
\midrule
Fiber coherence & Stronger sufficient condition & Connected components are realized transcript fibers & $\log |\mathcal T_{\mathrm{real}}|$ \\
\midrule
Meet-witness & Sufficient for transitivity and collapse & Base confusability collapses to the component cluster graph & $\log |\Pi_{\mathrm{cc}}|$ \\
\midrule
Transitivity & Sufficient collapse condition used here & $G_{\mathrm{conf}}=\mathrm{Cluster}(\Pi_{\mathrm{cc}})$ & $\log |\Pi_{\mathrm{cc}}|$ \\
\bottomrule
\end{tabularx}
\endgroup
\caption{Hierarchy of sufficient conditions for the cluster-graph equality route.}
\label{tab:equality-hierarchy}
\end{table}
